\chapter*{Appendix F: Sheaf Theory and Descent in RSVP}
\addcontentsline{toc}{chapter}{Appendix F: Sheaf Theory and Descent in RSVP}

\section{Presheaves}

\subsection{Definition}

Let $(M,\tau)$ be a topological space.

\begin{definition}[Presheaf]
A presheaf $\mathcal{F}$ assigns to each open set $U\subseteq M$ a set $\mathcal{F}(U)$ and to each inclusion $V\subseteq U$ a restriction map
\[
\rho^U_V : \mathcal{F}(U) \to \mathcal{F}(V)
\]
such that:
\[
\rho^U_U = \mathrm{id}, \qquad
\rho^V_W \circ \rho^U_V = \rho^U_W.
\]
\end{definition}

Typical examples include presheaves of:
scalar fields, vector fields, embeddings, entropy assignments, and Polyxan labels.

\section{Sheaves}

\subsection{Sheaf Axioms}

\begin{definition}[Sheaf]
A presheaf $\mathcal{F}$ is a sheaf if for every open cover $\{U_i\}$ of $U$:

\begin{itemize}
\item (Locality)  
If $s,t\in\mathcal{F}(U)$ and $s|_{U_i}=t|_{U_i}$ for all $i$, then $s=t$.

\item (Gluing)  
If $s_i \in \mathcal{F}(U_i)$ satisfy  
$s_i|_{U_i\cap U_j} = s_j|_{U_i\cap U_j}$,  
then there exists a unique  
$s \in \mathcal{F}(U)$ with $s|_{U_i}=s_i$.
\end{itemize}
\end{definition}

\section{Sheafification}

\subsection{Plus Construction}

Let $\mathcal{F}$ be a presheaf.

The sheafification $\mathcal{F}^+$ consists of equivalence classes of matching families on open covers.  
The canonical morphism  
\[
\eta : \mathcal{F} \to \mathcal{F}^+
\]
is universal among maps from $\mathcal{F}$ to sheaves.

\section{Stalks}

\subsection{Definition}

For $x\in M$:

\[
\mathcal{F}_x = \varinjlim_{x\in U} \mathcal{F}(U)
\]

is the stalk at $x$.

Stalks record all local behaviours relevant for RSVP curvature and Polyxan incidence.

\section{Sheaf of RSVP Fields}

\subsection{Scalar, Vector, and Entropy Fields}

Define:

\[
\mathcal{F}_\Phi(U) = \{ \Phi: U\to\mathbb{R} \},
\qquad
\mathcal{F}_{\mathbf{v}}(U) = \Gamma(TU),
\qquad
\mathcal{F}_S(U) = \{ S: U\to\mathbb{R} \}.
\]

These are sheaves under restriction because smooth fields satisfy locality and gluing.

\section{Sheaf of Embeddings}

\subsection{Incidence Embedding Sheaf}

For a fixed hypergraph $\mathcal{G}$:

\[
\mathcal{F}_{\iota}(U) =
\{ \iota : \mathrm{Inc}(\mathcal{G}|_U)\to U \ \text{smooth and fiber-preserving} \}.
\]

\begin{proposition}
$\mathcal{F}_\iota$ is a sheaf on $M$.
\end{proposition}

\section{Descent}

\subsection{Descent Datum}

Given a cover $\{U_i\}$, a descent datum consists of sections $s_i \in \mathcal{F}(U_i)$ and isomorphisms
\[
\phi_{ij} : s_i|_{U_i\cap U_j} \to s_j|_{U_i\cap U_j}
\]
satisfying the cocycle condition on triple overlaps.

\begin{definition}[Effective Descent]
A sheaf satisfies effective descent if every descent datum glues to a global section.
\end{definition}

\section{Hypergraph Sheaves}

\subsection{Polyxan Incidence Sheaf}

For each cell $c \in \mathcal{G}$, define a presheaf
\[
\mathcal{P}_c(U)
  = \{ \text{incidences of $c$ realisable in $U$} \}.
\]

The assignment
\[
\mathcal{P}(U)=\prod_{c\in \mathcal{G}} \mathcal{P}_c(U)
\]
is the Polyxan incidence sheaf.

\section{Curvature Sheaves}

\begin{definition}[Curvature Presheaf]
\[
\mathcal{R}(U) = \{ \mathcal{R}|_U : U\to\mathbb{R} \}.
\]
\end{definition}

\begin{proposition}
The curvature presheaf is a sheaf.
\end{proposition}

\section{Sheaf Cochains}

Let $\mathcal{F}$ be a sheaf of abelian groups.  
Define:

\[
C^k(\mathfrak{U}, \mathcal{F})
  = \prod_{i_0 < \dots < i_k}
      \mathcal{F}(U_{i_0}\cap\cdots\cap U_{i_k})
\]

with Čech differential.

\section{Čech Cohomology}

\subsection{Definition}

\[
\check{H}^k(M,\mathcal{F})
  = \frac{\ker(d: C^k\to C^{k+1})}
         {\mathrm{im}(d: C^{k-1}\to C^k)}.
\]

These groups detect obstructions to gluing—used in reconciliation and embedding stability.

\section{Sheaf-Theoretic Reconciliation}

\subsection{Local Reconciliation}

On $U$ and $V$ with overlap:
\[
p_U|_{U\cap V} = p_V|_{U\cap V}.
\]

\subsection{Global Reconciliation}

The global patch is the unique glued section:
\[
p = \mathrm{glue}\{p_U\}.
\]

\section{Descent for Semantic Patches}

\begin{definition}[Patch Sheaf]
\[
\mathcal{P}(U) = \{ p: \text{semantic patch supported in } U\}.
\]
\end{definition}

\begin{proposition}
$\mathcal{P}$ is a sheaf.
\end{proposition}

Thus all local semantic patches glue uniquely to a global patch.

\section{Sheaf-Theoretic Interpretation of Synthetic Duality}

The duality functor  
\[
\mathbb{D}: \mathbf{RSVP} \leftrightarrows \mathbf{TypHyper}
\]
is a morphism of sheaf topoi:
\[
\mathbb{D}^\ast : \mathbf{Sh}(M) \to \mathbf{Sh}(|\mathcal{G}|).
\]

\section{Sheaf Conditions for Normal Forms}

Let $\mathrm{nf}$ be semantic patch normalisation.

\begin{proposition}
$\mathrm{nf}$ is a sheaf morphism:
\[
\mathrm{nf}(p|_U) = (\mathrm{nf}(p))|_U.
\]
\end{proposition}

\section{Local–Global Consistency}

A configuration is globally curvature-aligned iff it is locally curvature-aligned on a good cover.

\section{Sheaf Models of Embedding Flows}

Embedding flows
\[
t\mapsto \iota_t
\]
form a sheaf of time-indexed sections over $M\times\mathbb{R}$.

\section{Semantic Descent Conditions}

For RSVP–Polyxan compatibility:

\[
\iota_U \simeq \iota_V \quad \text{on } U\cap V
\]

ensures globally well-defined embeddings.

\section{Higher Descent and Stacks}

For quine-stable embeddings, descent data naturally define a stack of curvature-preserving embeddings:
\[
\mathfrak{E}(U) = \{ \text{embeddings } \iota: \mathrm{Inc}(\mathcal{G})|_U \to U \}.
\]

\section{Stacky Patches}

Semantic patches form a fibered category over $M$; their stackification accounts for higher coherences in reconciliation.

\section{Obstruction Theory}

Obstructions to global embedding or global patch extension sit in:
\[
\check{H}^2(M,\mathcal{F}_\iota),
\qquad
\check{H}^1(M,\mathcal{P}).
\]

\section{Vanishing Obstructions Under Reconciliation}

Reconciliation eliminates Čech 1-cocycles:
\[
\check{H}^1(M,\mathcal{P}) = 0,
\]
so global patches always exist and are unique.


